% Source: https://artofproblemsolving.com/community/c3h1772052_middle_school_math_challenge_question_square_roots_and_exponents
% Stable identity: aops-topic-1772052-q01
Interpret the expression as the limit of $x_1=\sqrt2$, $x_{n+1}=\sqrt{2-x_n}$. All terms lie between $0$ and $\sqrt2$, and
\[
|x_{n+1}-1|=\left|\sqrt{2-x_n}-1\right|=\dfrac{|1-x_n|}{\sqrt{2-x_n}+1}\le|x_n-1| \cdot\dfrac{1}{1+\sqrt{2-\sqrt2}}<0.7\,|x_n-1| ,
\]
so the terms approach $1$. Equivalently, the value $x\ge0$ must satisfy $x=\sqrt{2-x}$, so $x^2+x-2=0$, $(x+2)(x-1)=0$, and since $x\ge0$, $x=\boxed{1}$.
